\section{信頼性と故障}

\subsection{平成27年度 問7}
\label{q:h27-reg-apt-07}

\noindent\textbf{原典：} sources/H27/aptitude.pdf、PDF物理7頁。

\textbf{公式正答：} 1（single）。

\par\textit{原典本文の校正と独自解説は未作成。分類のみの登録項目。}

\subsection{令和3年度 問6}
\label{q:r03-reg-apt-06}

\noindent\textbf{原典：} sources/R03/aptitude.pdf、PDF物理6頁。

\textbf{公式正答：} 1（single）。

\par\textit{原典本文の校正と独自解説は未作成。分類のみの登録項目。}

\section{工程管理・PERT}

\subsection{令和元年度 問10}
\label{q:r01-reg-apt-10}

\noindent\textbf{原典：} sources/R01/aptitude.pdf、PDF物理11頁。

\textbf{公式正答：} 3（single）。

\par\textit{原典本文の校正と独自解説は未作成。分類のみの登録項目。}

\section{最適化・意思決定}

\subsection{平成25年度 問14}
\label{q:h25-reg-apt-14}

\noindent\textbf{原典：} sources/H25/aptitude.pdf、PDF物理15頁。

\textbf{公式正答：} 1（single）。

\par\textit{原典本文の校正と独自解説は未作成。分類のみの登録項目。}

\subsection{平成26年度 問4}
\label{q:h26-reg-apt-04}

\noindent\textbf{原典：} sources/H26/aptitude.pdf、PDF物理4頁。

\textbf{公式正答：} 4（single）。

\par\textit{原典本文の校正と独自解説は未作成。分類のみの登録項目。}

\subsection{令和元年度 問4}
\label{q:r01-reg-apt-04}

\noindent\textbf{原典：} sources/R01/aptitude.pdf、PDF物理5頁。

\textbf{公式正答：} 1（single）。

\par\textit{原典本文の校正と独自解説は未作成。分類のみの登録項目。}

\subsection{令和3年度 問3}
\label{q:r03-reg-apt-03}

\noindent\textbf{原典：} sources/R03/aptitude.pdf、PDF物理3頁。

\textbf{公式正答：} 1（single）。

\par\textit{原典本文の校正と独自解説は未作成。分類のみの登録項目。}

\subsection{令和6年度 問9}
\label{q:r06-reg-apt-09}

\noindent\textbf{原典：} sources/R06/aptitude.pdf、PDF物理9頁。

\textbf{公式正答：} 5（single）。

\par\textit{原典本文の校正と独自解説は未作成。分類のみの登録項目。}

\section{暗号・認証・攻撃対策}

\subsection{平成28年度 問13}
\label{q:h28-reg-apt-13}

\noindent\textbf{原典：} sources/H28/aptitude.pdf、PDF物理13頁。

\textbf{公式正答：} 4（single）。

\par\textit{原典本文の校正と独自解説は未作成。分類のみの登録項目。}

\subsection{平成29年度 問6}
\label{q:h29-reg-apt-06}

\noindent\textbf{原典：} sources/H29/aptitude.pdf、PDF物理6頁。

\textbf{公式正答：} 1（single）。

\par\textit{原典本文の校正と独自解説は未作成。分類のみの登録項目。}

\subsection{令和4年度 問13}
\label{q:r04-reg-apt-13}

\noindent\textbf{原典：} sources/R04/aptitude.pdf、PDF物理12頁。

\textbf{公式正答：} 3（single）。

\par\textit{原典本文の校正と独自解説は未作成。分類のみの登録項目。}

\section{IP・通信}

\subsection{平成24年度 問4}
\label{q:h24-reg-apt-04}

\noindent\textbf{原典：} sources/H24/aptitude.pdf、PDF物理3頁。

\textbf{公式正答：} 2（single）。

\par\textit{原典本文の校正と独自解説は未作成。分類のみの登録項目。}

\section{アルゴリズム・再帰}

\subsection{令和3年度 問14}
\label{q:r03-reg-apt-14}

\noindent\textbf{原典：} sources/R03/aptitude.pdf、PDF物理14頁。

\textbf{公式正答：} 5（single）。

\par\textit{原典本文の校正と独自解説は未作成。分類のみの登録項目。}

\section{生物・遺伝}

\subsection{令和2年度 問14}
\label{q:r02-reg-apt-14}

\noindent\textbf{原典：} sources/R02/aptitude.pdf、PDF物理14頁。

\textbf{公式正答：} 1（single）。

\par\textit{原典本文の校正と独自解説は未作成。分類のみの登録項目。}

\section{環境法・循環・持続可能性}

\subsection{平成23年度 問2}
\label{q:h23-reg-apt-02}

\noindent\textbf{原典：} sources/H23/aptitude.pdf、PDF物理3頁。

\textbf{公式正答：} 3（single）。

\par\textit{原典本文の校正と独自解説は未作成。分類のみの登録項目。}

\subsection{令和2年度 問10}
\label{q:r02-reg-apt-10}

\noindent\textbf{原典：} sources/R02/aptitude.pdf、PDF物理10頁。

\textbf{公式正答：} 1（single）。

\par\textit{原典本文の校正と独自解説は未作成。分類のみの登録項目。}

\subsection{令和3年度 問4}
\label{q:r03-reg-apt-04}

\noindent\textbf{原典：} sources/R03/aptitude.pdf、PDF物理4頁。

\textbf{公式正答：} 2（single）。

\par\textit{原典本文の校正と独自解説は未作成。分類のみの登録項目。}

\subsection{令和6年度 問11}
\label{q:r06-reg-apt-11}

\noindent\textbf{原典：} sources/R06/aptitude.pdf、PDF物理11頁。

\textbf{公式正答：} 5（single）。

\par\textit{原典本文の校正と独自解説は未作成。分類のみの登録項目。}

\section{公益・説明責任}

\subsection{平成25年度 問2}
\label{q:h25-reg-apt-02}

\noindent\textbf{原典：} sources/H25/aptitude.pdf、PDF物理3頁。

\textbf{公式正答：} 3（single）。

\par\textit{原典本文の校正と独自解説は未作成。分類のみの登録項目。}

\section{安全・リスク評価}

\subsection{平成23年度 問3}
\label{q:h23-reg-apt-03}

\noindent\textbf{原典：} sources/H23/aptitude.pdf、PDF物理3頁。

\textbf{公式正答：} 5（single）。

\par\textit{原典本文の校正と独自解説は未作成。分類のみの登録項目。}

\subsection{平成23年度 問7}
\label{q:h23-reg-apt-07}

\noindent\textbf{原典：} sources/H23/aptitude.pdf、PDF物理8頁。

\textbf{公式正答：} 5（single）。

\par\textit{原典本文の校正と独自解説は未作成。分類のみの登録項目。}

\subsection{令和4年度 問4}
\label{q:r04-reg-apt-04}

\noindent\textbf{原典：} sources/R04/aptitude.pdf、PDF物理3頁。

\textbf{公式正答：} 1（single）。

\par\textit{原典本文の校正と独自解説は未作成。分類のみの登録項目。}

\subsection{令和7年度 問9}
\label{q:r07-reg-apt-09}

\noindent\textbf{原典：} sources/R07/aptitude.pdf、PDF物理9頁。

\textbf{公式正答：} 4（single）。

\par\textit{原典本文の校正と独自解説は未作成。分類のみの登録項目。}

\section{研究・情報倫理}

\subsection{平成24年度 問15}
\label{q:h24-reg-apt-15}

\noindent\textbf{原典：} sources/H24/aptitude.pdf、PDF物理13頁。

\textbf{公式正答：} 5（single）。

\par\textit{原典本文の校正と独自解説は未作成。分類のみの登録項目。}

\subsection{平成26年度 問12}
\label{q:h26-reg-apt-12}

\noindent\textbf{原典：} sources/H26/aptitude.pdf、PDF物理12頁。

\textbf{公式正答：} 2（single）。

\par\textit{原典本文の校正と独自解説は未作成。分類のみの登録項目。}

\subsection{平成26年度 問2}
\label{q:h26-reg-apt-02}

\noindent\textbf{原典：} sources/H26/aptitude.pdf、PDF物理2頁。

\textbf{公式正答：} 5（single）。

\paragraph{問題の要旨} 2013年改訂の科学者の行動規範から、最も不適切な記述を選ぶ。

\Needspace{4\baselineskip}\paragraph{確認した条件・記述（独自要約）}\mbox{}\par

\noindent\textbf{時点：} 2006年制定、2013年改訂の日本学術会議声明。\par

\noindent\textbf{方向：} 不適切な記述を1つ選ぶ。\par

\paragraph{問われる核心} 科学の自律性と社会への責任

\paragraph{必要知識} 科学者の範囲、研究の功績と責任、説明責任、社会との対話。

\paragraph{考え方・正答を選ぶ理由（査読待ち初稿）} 各肢の主語と責任を確認する。1〜4は科学者の範囲、品質責任、説明、研究公正という規範の要素に対応する。5は科学の自律性を社会からの孤立と読み替え、対話と助言を否定しているため選ぶ。

\Needspace{4\baselineskip}\paragraph{選択肢の検討（査読待ち初稿）}\mbox{}\par

\noindent\textbf{1：問の条件に一致しない候補} 科学者の範囲は所属に依存せず人文・社会・自然科学の研究者と知識を活用する専門職を含む。前文での科学者の定義に合う。自然科学者や大学所属者に限定しない。\par

\noindent\textbf{2：問の条件に一致しない候補} 知識・技術の質を保証し、健康・福祉・安全・地球環境へ貢献する責任を持つ。基本的責任の規定に合う。専門知識の質と社会への貢献の双方を求める。\par

\noindent\textbf{3：問の条件に一致しない候補} 研究の意義と社会・環境への影響を中立・客観的に説明・公表する。説明と公開の責務に合う。専門性を理由に説明責任を免れるものではない。\par

\noindent\textbf{4：問の条件に一致しない候補} 役割に応じた功績と責任を負い、データを保存してFFPを行わず加担しない。研究活動の公正さに関する責務に合う。功績のみを受け取る立場ではない。\par

\noindent\textbf{5：問の条件に一致する候補} 独立性のため市民対話に否定的であり、政策への科学的助言も避ける。独立した科学的判断は対話や助言を避ける理由にならない。規範は市民との交流と政策への助言を促すので不適切である。\par

\paragraph{再利用する考え方} 制度の対象・定義と倫理上の責任を区別し、科学の自律性と社会への責任を具体的行為に適用する。

\paragraph{共通理論} \hyperref[theory:research:norms]{第\ref*{theory:research:norms}節}

\par 参考：\cite{scj_code2013}。

\par 難易度 2（校準初稿）、優先度未評価。

\subsection{平成26年度 問3}
\label{q:h26-reg-apt-03}

\noindent\textbf{原典：} sources/H26/aptitude.pdf、PDF物理3頁。

\textbf{公式正答：} 1（single）。

\paragraph{問題の要旨} 捏造・改ざん・盗用という3種の研究不正の定義を判定する。

\Needspace{4\baselineskip}\paragraph{確認した条件・記述（独自要約）}\mbox{}\par

\noindent\textbf{判断順：} 記述アから順に○=適切、×=不適切。\par

\noindent\textbf{ア：} 捏造は存在しないデータや研究結果を作ること。\par

\noindent\textbf{イ：} 改ざんは資料・機器・過程などを操作し、データや結果を真正でないものに加工すること。\par

\noindent\textbf{ウ：} 盗用は他者のアイデア・分析法・データ・結果・論文・用語を、了解若しくは適切な表示なしに流用すること。\par

\paragraph{問われる核心} 研究不正FFPの定義

\paragraph{必要知識} 捏造・改ざん・盗用の対象、データ処理と不真正化の区別。

\paragraph{考え方・正答を選ぶ理由（査読待ち初稿）} アは架空のデータの作成、イは取得・処理過程を操作した不真正化、ウは他者の成果の不適切な流用という各定義に合う。3記述とも○なので第1候補である。引用を示す責任と、著作物利用の許諾の問題は別に検討し、単に許諾があれば他者の成果を自分のものと表示してよいとは考えない。

\Needspace{4\baselineskip}\paragraph{選択肢の検討（査読待ち初稿）}\mbox{}\par

\noindent\textbf{1：問の条件に一致する候補} ○○○。全記述の判定○○○に一致する。\par

\noindent\textbf{2：問の条件に一致しない候補} ×○○。記述アの判定が逆である。独立判定は○○○となる。\par

\noindent\textbf{3：問の条件に一致しない候補} ○×○。記述イの判定が逆である。独立判定は○○○となる。\par

\noindent\textbf{4：問の条件に一致しない候補} ××○。記述ア・イの判定が逆である。独立判定は○○○となる。\par

\noindent\textbf{5：問の条件に一致しない候補} ×××。記述ア・イ・ウの判定が逆である。独立判定は○○○となる。\par

\paragraph{再利用する考え方} 制度の対象・定義と倫理上の責任を区別し、研究不正FFPの定義を具体的行為に適用する。

\paragraph{共通理論} \hyperref[theory:research:ffp]{第\ref*{theory:research:ffp}節}

\par 参考：\cite{mext_misconduct2014}。

\par 難易度 2（校準初稿）、優先度未評価。

\subsection{平成27年度 問8}
\label{q:h27-reg-apt-08}

\noindent\textbf{原典：} sources/H27/aptitude.pdf、PDF物理8頁。

\textbf{公式正答：} 1（single）。

\paragraph{問題の要旨} 実験ノートの記録・保存・訂正・異動時の持出しを判定する。

\Needspace{4\baselineskip}\paragraph{確認した条件・記述（独自要約）}\mbox{}\par

\noindent\textbf{判断順：} 記述アから順に○=適切、×=不適切。\par

\noindent\textbf{ア：} 原則として綴じたノートを消せない筆記具で記載する。\par

\noindent\textbf{イ：} 研究ノートは特許紛争回避や権利関係の明確化にも役立ち、企業研究者にも有用である。\par

\noindent\textbf{ウ：} 研究機関は一定期間の研究データ保存と必要時の開示を義務付ける規程を整える。\par

\noindent\textbf{エ：} 記述の誤りがあるページは廃棄し、見やすく修正するのが望ましい。\par

\noindent\textbf{オ：} 研究者個人に帰属するので異動時には個人判断で持ち出せる。\par

\paragraph{問われる核心} 研究記録の追跡可能性と機関管理

\paragraph{必要知識} 元記録の保存、訂正の履歴、データ帰属と異動時の管理。

\paragraph{考え方・正答を選ぶ理由（査読待ち初稿）} アは紙ノートの追記・差替えを追跡できるようにする原則である。イは研究の経緯を立証する記録として有用という記述であり、ノートだけで特許権が自動的に決まるとの意味ではない。ウは機関の保存規程に対応する。エのページ廃棄は元の経過を失わせる。オの個人所有・自由持出しの一律断定も不適切であり、資金・契約・機関規程による管理が必要である。○○○××の第1候補を選ぶ。

\Needspace{4\baselineskip}\paragraph{選択肢の検討（査読待ち初稿）}\mbox{}\par

\noindent\textbf{1：問の条件に一致する候補} ○○○××。全記述の判定○○○××に一致する。\par

\noindent\textbf{2：問の条件に一致しない候補} ×○○××。記述アの判定が逆である。独立判定は○○○××となる。\par

\noindent\textbf{3：問の条件に一致しない候補} ○×○××。記述イの判定が逆である。独立判定は○○○××となる。\par

\noindent\textbf{4：問の条件に一致しない候補} ×○×○×。記述ア・ウ・エの判定が逆である。独立判定は○○○××となる。\par

\noindent\textbf{5：問の条件に一致しない候補} ○○○×○。記述オの判定が逆である。独立判定は○○○××となる。\par

\paragraph{再利用する考え方} 制度の対象・定義と倫理上の責任を区別し、研究記録の追跡可能性と機関管理を具体的行為に適用する。

\paragraph{共通理論} \hyperref[theory:research:records]{第\ref*{theory:research:records}節}

\par 参考：\cite{mext_misconduct2014,jsps_integrity2025}。

\par 難易度 2（校準初稿）、優先度未評価。

\subsection{平成28年度 問4}
\label{q:h28-reg-apt-04}

\noindent\textbf{原典：} sources/H28/aptitude.pdf、PDF物理4頁。

\textbf{公式正答：} 2（single）。

\paragraph{問題の要旨} 2014年ガイドラインの特定不正行為と研究倫理全般を区別して判定する。

\Needspace{4\baselineskip}\paragraph{確認した条件・記述（独自要約）}\mbox{}\par

\noindent\textbf{判断順：} 記述アから順に○=適切、×=不適切。\par

\noindent\textbf{ア：} 故意又は基本的注意義務を著しく怠った、発表成果中の捏造・改ざん・盗用が特定不正行為である。\par

\noindent\textbf{イ：} 確認された不正に係る競争的資金について配分機関は事案に応じ交付取消しや一部・全部の返還を求められる。\par

\noindent\textbf{ウ：} 二重投稿や不適切な著者表示は研究者倫理に反するが不正行為ではない。\par

\noindent\textbf{エ：} 不正の対応は研究者・コミュニティ・機関の規律と自浄作用に基づく。\par

\noindent\textbf{オ：} 研究倫理教育は倫理を向上させるが、不正予防の環境整備にはならない。\par

\paragraph{問われる核心} 特定不正行為と広い研究倫理違反

\paragraph{必要知識} FFPの限定、一般的な不正、競争的資金の措置、研究倫理教育。

\paragraph{考え方・正答を選ぶ理由（査読待ち初稿）} アは特定不正行為の要件である。イは対象の競争的資金と事案に応じた措置に限る記述で妥当であり、全資金の自動返還ではない。ウは特定不正行為の限定を不正一般へ広げている。二重投稿等も倫理に反する不正として認識されるので×である。エは自己規律・自律の位置付けに合い、オは教育の予防的役割を否定するため×である。○○×○×の第2候補となる。

\Needspace{4\baselineskip}\paragraph{選択肢の検討（査読待ち初稿）}\mbox{}\par

\noindent\textbf{1：問の条件に一致しない候補} ○○○○×。記述ウの判定が逆である。独立判定は○○×○×となる。\par

\noindent\textbf{2：問の条件に一致する候補} ○○×○×。全記述の判定○○×○×に一致する。\par

\noindent\textbf{3：問の条件に一致しない候補} ○×○×○。記述イ・ウ・エ・オの判定が逆である。独立判定は○○×○×となる。\par

\noindent\textbf{4：問の条件に一致しない候補} ×○×○×。記述アの判定が逆である。独立判定は○○×○×となる。\par

\noindent\textbf{5：問の条件に一致しない候補} ○×○××。記述イ・ウ・エの判定が逆である。独立判定は○○×○×となる。\par

\paragraph{再利用する考え方} 制度の対象・定義と倫理上の責任を区別し、特定不正行為と広い研究倫理違反を具体的行為に適用する。

\paragraph{共通理論} \hyperref[theory:research:ffp]{第\ref*{theory:research:ffp}節}, \hyperref[theory:research:prevention]{第\ref*{theory:research:prevention}節}

\par 参考：\cite{mext_misconduct2014}。

\par 難易度 2（校準初稿）、優先度未評価。

\subsection{平成29年度 問14}
\label{q:h29-reg-apt-14}

\noindent\textbf{原典：} sources/H29/aptitude.pdf、PDF物理13頁。

\textbf{公式正答：} 3（single）。

\paragraph{問題の要旨} 研究の自由、記録、資金、著者責任、引用に関する7記述中の不適切数を求める。

\Needspace{4\baselineskip}\paragraph{確認した条件・記述（独自要約）}\mbox{}\par

\noindent\textbf{方向：} 不適切な記述の個数。候補1〜5はそれぞれ0〜4個。\par

\noindent\textbf{ア：} 研究の自由は真理又は公益、独創性と正確さ、結果への責任を伴う。\par

\noindent\textbf{イ：} 客観的検証のため記録を保存し、FFPを禁じる。\par

\noindent\textbf{ウ：} 研究費は計画・申請に基づき公正に使い、目的外利用・不正操作をしない。\par

\noindent\textbf{エ：} 貢献者を著者とし論文の正確さなどへ責任を負うが、共著者は自分に関係した部分だけへ責任を持てばよい。\par

\noindent\textbf{オ：} 実験の大きな貢献があれば、論文の内容と正確さを説明できなくても共著者になれる。\par

\noindent\textbf{カ：} 先行研究を明示し、自らの独創的寄与を区別する。\par

\noindent\textbf{キ：} 著作物の引用部分と原著者を明示し、引用の範囲を超えるコピーには原則として許諾を得る。\par

\paragraph{問われる核心} 著者の貢献と説明責任

\paragraph{必要知識} 功績の帰属、共同研究での責任、引用と著作物利用の区別。

\paragraph{考え方・正答を選ぶ理由（査読待ち初稿）} ア・イ・ウは研究の目的、検証、資金責任、カ・キは先行成果と利用条件の明示で適切である。エは一律に自分の部分だけへ責任を限定しており不適切である。役割分担を明示する場合も、共同研究の信頼性の確認を放棄することにはならない。オは実験への貢献だけで内容を説明できない人を著者にできるとする点が不適切である。誤りは2件、第3候補となる。

\Needspace{4\baselineskip}\paragraph{選択肢の検討（査読待ち初稿）}\mbox{}\par

\noindent\textbf{1：問の条件に一致しない候補} 0個。エ・オの2つの誤りを数えていない。\par

\noindent\textbf{2：問の条件に一致しない候補} 1個。エの責任限定とオの説明不能の2つが不適切である。\par

\noindent\textbf{3：問の条件に一致する候補} 2個。エ・オの2記述を不適切と数えた結果に一致する。\par

\noindent\textbf{4：問の条件に一致しない候補} 3個。ア・イ・ウ・カ・キには、記載された倫理原則に反する内容はない。\par

\noindent\textbf{5：問の条件に一致しない候補} 4個。不適切はエ・オだけであり、4件ではない。\par

\paragraph{再利用する考え方} 制度の対象・定義と倫理上の責任を区別し、著者の貢献と説明責任を具体的行為に適用する。

\paragraph{共通理論} \hyperref[theory:research:authorship]{第\ref*{theory:research:authorship}節}

\par 参考：\cite{jsps_integrity2025}。

\par 難易度 2（校準初稿）、優先度未評価。

\subsection{令和元年度再試験 問10}
\label{q:r01-ret-apt-10}

\noindent\textbf{原典：} sources/R01\_retest/aptitude.pdf、PDF物理10頁。

\textbf{公式正答：} 2（single）。

\paragraph{問題の要旨} 誠実な研究の誤り、私費研究、言語の異なる重複発表、共著論文の翻訳利用を判定する。

\Needspace{4\baselineskip}\paragraph{確認した条件・記述（独自要約）}\mbox{}\par

\noindent\textbf{判断順：} 記述アから順に○=適切、×=不適切。\par

\noindent\textbf{ア：} 適切な方法で正当に得た研究成果が結果的に誤りでも、ガイドライン公表後は不正行為となった。\par

\noindent\textbf{イ：} 個人の資金と労力で研究・公表するならガイドラインの内容を考慮する必要がない。\par

\noindent\textbf{ウ：} 同一成果を日本語と英語で別の学会に発表すれば二重投稿にはならない。\par

\noindent\textbf{エ：} AとBの共著英語論文の一部を、Aが単独日本語書籍に翻訳利用する前にBへ相談して了解を得た。\par

\paragraph{問われる核心} 研究の誠実さと公表時の透明性

\paragraph{必要知識} 誠実な誤りの除外、資金源に依存しない倫理、重複発表と共同成果利用。

\paragraph{考え方・正答を選ぶ理由（査読待ち初稿）} アは誠実な誤りを不正とする点、イは資金源だけで倫理責任を免れる点、ウは言語だけで重複発表を正当化する点がいずれも誤りである。エは共同成果の再利用を独断で行わず共著者へ相談した行為として適切である。出版側の許諾・引用表示などが不要との記述ではなく、実利用時には別途確認する。×××○の第2候補である。

\Needspace{4\baselineskip}\paragraph{選択肢の検討（査読待ち初稿）}\mbox{}\par

\noindent\textbf{1：問の条件に一致しない候補} ×○×○。記述イの判定が逆である。独立判定は×××○となる。\par

\noindent\textbf{2：問の条件に一致する候補} ×××○。全記述の判定×××○に一致する。\par

\noindent\textbf{3：問の条件に一致しない候補} ○××○。記述アの判定が逆である。独立判定は×××○となる。\par

\noindent\textbf{4：問の条件に一致しない候補} ○○○×。記述ア・イ・ウ・エの判定が逆である。独立判定は×××○となる。\par

\noindent\textbf{5：問の条件に一致しない候補} ××○○。記述ウの判定が逆である。独立判定は×××○となる。\par

\paragraph{再利用する考え方} 制度の対象・定義と倫理上の責任を区別し、研究の誠実さと公表時の透明性を具体的行為に適用する。

\paragraph{共通理論} \hyperref[theory:research:ffp]{第\ref*{theory:research:ffp}節}, \hyperref[theory:research:authorship]{第\ref*{theory:research:authorship}節}

\par 参考：\cite{mext_misconduct2014,jsps_integrity2025}。

\par 難易度 2（校準初稿）、優先度未評価。

\subsection{令和元年度再試験 問15}
\label{q:r01-ret-apt-15}

\noindent\textbf{原典：} sources/R01\_retest/aptitude.pdf、PDF物理16頁。

\textbf{公式正答：} 1（single）。

\par\textit{原典本文の校正と独自解説は未作成。分類のみの登録項目。}

\subsection{令和元年度再試験 問6}
\label{q:r01-ret-apt-06}

\noindent\textbf{原典：} sources/R01\_retest/aptitude.pdf、PDF物理6頁。

\textbf{公式正答：} 4（single）。

\paragraph{問題の要旨} 科学者の範囲・誠実さ・両義的利用・研究環境の責任を判定する。

\Needspace{4\baselineskip}\paragraph{確認した条件・記述（独自要約）}\mbox{}\par

\noindent\textbf{判断順：} 記述アから順に○=適切、×=不適切。\par

\noindent\textbf{ア：} 科学者は所属にかかわらず全学術分野の研究者や科学知識の利活用を行う専門職を含む。\par

\noindent\textbf{イ：} 正直・誠実に行動し、専門能力の維持向上と知識の正確さの提示に努める。\par

\noindent\textbf{ウ：} 意図に反し悪用され得ても、社会へ貢献すると判断すれば速やかに研究・公表を積極的に行う。\par

\noindent\textbf{エ：} 公正な研究環境の維持と不正抑止の教育啓発に継続して取り組む責任がある。\par

\paragraph{問われる核心} 科学研究の利用の両義性

\paragraph{必要知識} 研究の便益と悪用リスク、社会に許容される研究・公表方法。

\paragraph{考え方・正答を選ぶ理由（査読待ち初稿）} アは前文の定義、イは誠実な姿勢、エは環境整備と教育の責任に対応する。ウは善用の可能性だけで悪用への配慮を省いている。規範は研究と公表の際に社会へ許容される手段・方法を選ぶことを求めるので×である。○○×○の第4候補である。

\Needspace{4\baselineskip}\paragraph{選択肢の検討（査読待ち初稿）}\mbox{}\par

\noindent\textbf{1：問の条件に一致しない候補} ○○○○。記述ウの判定が逆である。独立判定は○○×○となる。\par

\noindent\textbf{2：問の条件に一致しない候補} ×○○○。記述ア・ウの判定が逆である。独立判定は○○×○となる。\par

\noindent\textbf{3：問の条件に一致しない候補} ○×○○。記述イ・ウの判定が逆である。独立判定は○○×○となる。\par

\noindent\textbf{4：問の条件に一致する候補} ○○×○。全記述の判定○○×○に一致する。\par

\noindent\textbf{5：問の条件に一致しない候補} ○○○×。記述ウ・エの判定が逆である。独立判定は○○×○となる。\par

\paragraph{再利用する考え方} 制度の対象・定義と倫理上の責任を区別し、科学研究の利用の両義性を具体的行為に適用する。

\paragraph{共通理論} \hyperref[theory:research:norms]{第\ref*{theory:research:norms}節}, \hyperref[theory:research:prevention]{第\ref*{theory:research:prevention}節}

\par 参考：\cite{scj_code2013}。

\par 難易度 2（校準初稿）、優先度未評価。

\subsection{令和2年度 問3}
\label{q:r02-reg-apt-03}

\noindent\textbf{原典：} sources/R02/aptitude.pdf、PDF物理3頁。

\textbf{公式正答：} 3（single）。

\paragraph{問題の要旨} 兼業変更、競争相手の査読、配偶者名義の株式、公的設備の企業利用の利益相反管理を判定する。

\Needspace{4\baselineskip}\paragraph{確認した条件・記述（独自要約）}\mbox{}\par

\noindent\textbf{判断順：} 記述アから順に○=適切、×=不適切。\par

\noindent\textbf{ア：} 水曜日に許可された有償兼業について土曜の技術会議参加が必要となり、大学COI委員会へ相談した。\par

\noindent\textbf{イ：} 自己の研究と競争関係にある論文の公正な査読に不安があるため辞退した。\par

\noindent\textbf{ウ：} 公的資金で評価した企業の未公開株を少額で配偶者名義にて買い、COI委員会への相談を省いた。\par

\noindent\textbf{エ：} 大学の兼業許可を得て設立した企業が公的資金の研究設備を無償利用する必要が生じ、事前に大学COI委員会へ相談した。\par

\paragraph{問われる核心} 利益相反の開示と第三者管理

\paragraph{必要知識} 経済的利益と責務の衝突、家族の利害、査読の競争関係、機関資源。

\paragraph{考え方・正答を選ぶ理由（査読待ち初稿）} アは許可条件の変更を自己判断せず相談する対応、イは査読の公正さを確保する辞退で適切である。ウは家族名義や少額を理由に利害関係を隠しており不適切である。エは所属機関の資源と私企業の利益を分ける事前相談として適切である。相談すれば無償利用が必ず認められるとの意味ではない。○○×○の第3候補を選ぶ。

\Needspace{4\baselineskip}\paragraph{選択肢の検討（査読待ち初稿）}\mbox{}\par

\noindent\textbf{1：問の条件に一致しない候補} ○○○○。記述ウの判定が逆である。独立判定は○○×○となる。\par

\noindent\textbf{2：問の条件に一致しない候補} ○○○×。記述ウ・エの判定が逆である。独立判定は○○×○となる。\par

\noindent\textbf{3：問の条件に一致する候補} ○○×○。全記述の判定○○×○に一致する。\par

\noindent\textbf{4：問の条件に一致しない候補} ○×○○。記述イ・ウの判定が逆である。独立判定は○○×○となる。\par

\noindent\textbf{5：問の条件に一致しない候補} ×○○○。記述ア・ウの判定が逆である。独立判定は○○×○となる。\par

\paragraph{再利用する考え方} 制度の対象・定義と倫理上の責任を区別し、利益相反の開示と第三者管理を具体的行為に適用する。

\paragraph{共通理論} \hyperref[theory:research:coi]{第\ref*{theory:research:coi}節}

\par 参考：\cite{jsps_integrity2025}。

\par 難易度 2（校準初稿）、優先度未評価。

\subsection{令和2年度 問8}
\label{q:r02-reg-apt-08}

\noindent\textbf{原典：} sources/R02/aptitude.pdf、PDF物理8頁。

\textbf{公式正答：} 1（single）。

\par\textit{原典本文の校正と独自解説は未作成。分類のみの登録項目。}

\subsection{令和4年度 問15}
\label{q:r04-reg-apt-15}

\noindent\textbf{原典：} sources/R04/aptitude.pdf、PDF物理14頁。

\textbf{公式正答：} 5（single）。

\par\textit{原典本文の校正と独自解説は未作成。分類のみの登録項目。}

\subsection{令和4年度 問2}
\label{q:r04-reg-apt-02}

\noindent\textbf{原典：} sources/R04/aptitude.pdf、PDF物理2頁。

\textbf{公式正答：} 4（single）。

\par\textit{原典本文の校正と独自解説は未作成。分類のみの登録項目。}

\subsection{令和5年度 問7}
\label{q:r05-reg-apt-07}

\noindent\textbf{原典：} sources/R05/aptitude.pdf、PDF物理7頁。

\textbf{公式正答：} 2（single）。

\paragraph{問題の要旨} 科学者の研究公正、対話、科学的助言と政策判断の関係を判定する。

\Needspace{4\baselineskip}\paragraph{確認した条件・記述（独自要約）}\mbox{}\par

\noindent\textbf{判断順：} 記述アから順に○=適切、×=不適切。\par

\noindent\textbf{ア：} 公表は役割に応じた功績と責任を伴い、データを保存してFFPを行わず加担しない。\par

\noindent\textbf{イ：} 市民対話と政策助言に努め、科学者に意見の相違がある場合は多数決で統一見解を決めて助言する。\par

\noindent\textbf{ウ：} 客観的根拠に基づき権威を濫用せず、科学的不確実性と見解の多様性も説明する。\par

\noindent\textbf{エ：} 科学的知見は政策の唯一の根拠ではなく、助言と異なる決定には必要に応じ社会への説明を求める。\par

\paragraph{問われる核心} 科学的助言の不確実性と政策責任

\paragraph{必要知識} 合意と意見相違の説明、科学的根拠と政策決定の他の判断材料。

\paragraph{考え方・正答を選ぶ理由（査読待ち初稿）} ア・ウ・エはそれぞれ研究公正、助言の質、科学と政策の役割に合う。イの前半は適切だが、多数決で相違を消すという後半が誤りである。規範は意見が異なる場合に相違を説明することを求めている。○×○○の第2候補である。

\Needspace{4\baselineskip}\paragraph{選択肢の検討（査読待ち初稿）}\mbox{}\par

\noindent\textbf{1：問の条件に一致しない候補} ×○○○。記述ア・イの判定が逆である。独立判定は○×○○となる。\par

\noindent\textbf{2：問の条件に一致する候補} ○×○○。全記述の判定○×○○に一致する。\par

\noindent\textbf{3：問の条件に一致しない候補} ○○×○。記述イ・ウの判定が逆である。独立判定は○×○○となる。\par

\noindent\textbf{4：問の条件に一致しない候補} ○○○×。記述イ・エの判定が逆である。独立判定は○×○○となる。\par

\noindent\textbf{5：問の条件に一致しない候補} ○○○○。記述イの判定が逆である。独立判定は○×○○となる。\par

\paragraph{再利用する考え方} 制度の対象・定義と倫理上の責任を区別し、科学的助言の不確実性と政策責任を具体的行為に適用する。

\paragraph{共通理論} \hyperref[theory:research:norms]{第\ref*{theory:research:norms}節}

\par 参考：\cite{scj_code2013}。

\par 難易度 2（校準初稿）、優先度未評価。

\subsection{令和6年度 問8}
\label{q:r06-reg-apt-08}

\noindent\textbf{原典：} sources/R06/aptitude.pdf、PDF物理8頁。

\textbf{公式正答：} 3（single）。

\paragraph{問題の要旨} 贈与・幽霊著者、資金獲得のみの著者資格、細切れ出版を判定する。

\Needspace{4\baselineskip}\paragraph{確認した条件・記述（独自要約）}\mbox{}\par

\noindent\textbf{判断順：} 記述アから順に○=適切、×=不適切。\par

\noindent\textbf{ア：} 資格を満たさない者を好意や強制により著者へ入れる行為はギフト・オーサーシップという倫理違反である。\par

\noindent\textbf{イ：} 研究費を獲得した者は研究遂行へ寄与したので論文著者の資格を満たす。\par

\noindent\textbf{ウ：} 1つの研究を複数の小研究に分け細切れに公表することをサラミ又はボローニャ出版という。\par

\noindent\textbf{エ：} 実質的貢献と著者要件を満たす者を著者として表示しない行為はゴースト・オーサーシップという倫理違反である。\par

\paragraph{問われる核心} 著者資格と研究成果の適切な単位

\paragraph{必要知識} 貢献と責任、謝辞、ギフト・ゴースト著者、細切れ出版。

\paragraph{考え方・正答を選ぶ理由（査読待ち初稿）} アは不相応な著者の追加、エは資格ある著者の省略という定義に合う。イは資金獲得のみで著者資格を充足するとする誤りである。ウは細切れ公表の用語を説明しており○である。独立した科学的問いを持つ複数論文すべてが倫理違反という意味に一般化しない。○×○○の第3候補を選ぶ。

\Needspace{4\baselineskip}\paragraph{選択肢の検討（査読待ち初稿）}\mbox{}\par

\noindent\textbf{1：問の条件に一致しない候補} ○○○○。記述イの判定が逆である。独立判定は○×○○となる。\par

\noindent\textbf{2：問の条件に一致しない候補} ×○○○。記述ア・イの判定が逆である。独立判定は○×○○となる。\par

\noindent\textbf{3：問の条件に一致する候補} ○×○○。全記述の判定○×○○に一致する。\par

\noindent\textbf{4：問の条件に一致しない候補} ○○×○。記述イ・ウの判定が逆である。独立判定は○×○○となる。\par

\noindent\textbf{5：問の条件に一致しない候補} ○○○×。記述イ・エの判定が逆である。独立判定は○×○○となる。\par

\paragraph{再利用する考え方} 制度の対象・定義と倫理上の責任を区別し、著者資格と研究成果の適切な単位を具体的行為に適用する。

\paragraph{共通理論} \hyperref[theory:research:authorship]{第\ref*{theory:research:authorship}節}

\par 参考：\cite{jsps_integrity2025}。

\par 難易度 2（校準初稿）、優先度未評価。

\section{秘密情報の接近制御・相互確認と持出し防止}

\subsection{令和元年度 問9}
\label{q:r01-reg-apt-09}

\noindent\textbf{原典：} sources/R01/aptitude.pdf、PDF物理10頁。

\textbf{公式正答：} 2（single）。

\paragraph{問題の要旨} 秘密情報のアクセス、相互確認、表示、配置、暗号・閲覧制限、紙削減の6記述で不適切な数を求める。

\Needspace{4\baselineskip}\paragraph{確認した条件・記述（独自要約）}\mbox{}\par

\noindent\textbf{方向：} 不適切な記述の数。候補0/1/2/3/4。\par

\noindent\textbf{構成：} ア〜カの6項目。\par

\noindent\textbf{限定：} オは外部では閲覧できないように設定する対策。暗号自体の万能性や外部送信の許可ではない。\par

\paragraph{問われる核心} 秘密情報の接近制御と相互確認

\paragraph{必要知識} 秘密情報の接近制御と相互確認。ア○：情報の分類に応じたアクセス権で、必要のない人の接近を制限する。イ×：相互の確認や責任者のチェックを減らす単独作業の一律推奨は逆方向である。ウ○：秘密の表示で取扱者の認識を支える。エ○：資料・作業が死角になることを減らすが、権限のない者へ画面を見せてよい意味ではない。オ○：暗号と登録ID・端末等の閲覧権制限を適切に設定する対策である。暗号化だけで漏えいが絶対不可能になるという主張ではなく、無断送信を許すものでもない。カ○：電子的アクセス管理と組み合わせ、無権限者が紙に接する機会を減らす。不適切はイの1項目で2。

\paragraph{考え方・正答を選ぶ理由（査読待ち初稿）} ア○：情報の分類に応じたアクセス権で、必要のない人の接近を制限する。イ×：相互の確認や責任者のチェックを減らす単独作業の一律推奨は逆方向である。ウ○：秘密の表示で取扱者の認識を支える。エ○：資料・作業が死角になることを減らすが、権限のない者へ画面を見せてよい意味ではない。オ○：暗号と登録ID・端末等の閲覧権制限を適切に設定する対策である。暗号化だけで漏えいが絶対不可能になるという主張ではなく、無断送信を許すものでもない。カ○：電子的アクセス管理と組み合わせ、無権限者が紙に接する機会を減らす。不適切はイの1項目で2。

\Needspace{4\baselineskip}\paragraph{選択肢の検討（査読待ち初稿）}\mbox{}\par

\noindent\textbf{1：選択対象ではない} 0項目。単独作業を可能な限り推奨するイが不適切なので0ではない。\par

\noindent\textbf{2：選択対象} 1項目。イだけが不適切であり1項目となる。\par

\noindent\textbf{3：選択対象ではない} 2項目。イに加えて他の適切な対策も否定すると数え過ぎとなる。\par

\noindent\textbf{4：選択対象ではない} 3項目。アクセス権の設定・表示・配置・適切な閲覧制限・紙削減の全てを危険低減策として区別すれば3ではない。\par

\noindent\textbf{5：選択対象ではない} 4項目。適切な対策まで不適切に数えている。秘密保護と一人作業の推奨を混同しない。\par

\paragraph{再利用する考え方} 対象、条件、時点を先に確認する。秘密情報の接近制御と相互確認を具体的な事例へ適用し、技術的原因・法的要件・倫理的判断のどれを問うか区別する。

\paragraph{共通理論} \hyperref[theory:research:secret_management]{第\ref*{theory:research:secret_management}節}, \hyperref[theory:compliance:secret]{第\ref*{theory:compliance:secret}節}

\par 参考：\cite{meti_secret_handbook2024}。

\paragraph{関連問題} \hyperref[q:r05-reg-apt-02]{R05-REG-APT-02}

\par 難易度 2（校準初稿）、優先度未評価。

\subsection{令和5年度 問2}
\label{q:r05-reg-apt-02}

\noindent\textbf{原典：} sources/R05/aptitude.pdf、PDF物理2頁。

\textbf{公式正答：} 3（single）。

\paragraph{問題の要旨} アクセス権、秘密表示、単独作業、電子的コピー制限、座席配置という漏えい対策を判定する。

\Needspace{4\baselineskip}\paragraph{確認した条件・記述（独自要約）}\mbox{}\par

\noindent\textbf{方向：} 最も不適切。\par

\noindent\textbf{時点：} 令和5年。2024年版ハンドブックは現在の背景資料として区別する。\par

\noindent\textbf{範囲：} 5つの対策案。電子的制限の完全防御を主張していない。\par

\paragraph{問われる核心} 秘密保護におけるアクセス権と相互確認

\paragraph{必要知識} 秘密保護におけるアクセス権と相互確認。情報の接近制御、秘密の認識、持出しの困難化、相互チェックを組み合わせる。3だけはチェックを減らす単独作業を一律に推奨しているので不適切。複数人にアクセス権を無制限に与えることとは区別する。

\paragraph{考え方・正答を選ぶ理由（査読待ち初稿）} 情報の接近制御、秘密の認識、持出しの困難化、相互チェックを組み合わせる。3だけはチェックを減らす単独作業を一律に推奨しているので不適切。複数人にアクセス権を無制限に与えることとは区別する。

\Needspace{4\baselineskip}\paragraph{選択肢の検討（査読待ち初稿）}\mbox{}\par

\noindent\textbf{1：適切} 情報分類ごとにアクセス権の設定ルールと適切な範囲を定める。。必要性に基づくアクセス範囲の制御である。権限を無制限に広げる対策ではない。\par

\noindent\textbf{2：適切} 媒体へ秘密情報であることが分かる表示をする。。取扱者が秘密の分類と注意を認識する仕組みを支える。\par

\noindent\textbf{3：不適切} 複数人を避け、可能な限り単独で秘密情報を扱う。。必要に応じ複数人で相互確認できる体制を整え、単独が必要なら責任者が事前・事後に確認する。単独を最大化する方針は逆である。\par

\noindent\textbf{4：適切} 印刷・複製・外部メディアへのコピー等を制限する設定をする。。持出しの経路を減らす技術的対策となる。印刷禁止やコピーガードが撮影・再入力まで完全に防ぐとの保証ではない。\par

\noindent\textbf{5：適切} 作業・資料が死角になりにくい配置を検討する。。相互に作業へ注意が向く環境を作る。秘密画面をアクセス権のない人にさらすことを避けるという別条件も守る。\par

\paragraph{再利用する考え方} 対象、条件、時点を先に確認する。秘密保護におけるアクセス権と相互確認を具体的な事例へ適用し、技術的原因・法的要件・倫理的判断のどれを問うか区別する。

\paragraph{共通理論} \hyperref[theory:research:secret_management]{第\ref*{theory:research:secret_management}節}, \hyperref[theory:compliance:secret]{第\ref*{theory:compliance:secret}節}

\par 参考：\cite{meti_secret_handbook2024}。

\paragraph{関連問題} \hyperref[q:r01-reg-apt-09]{R01-REG-APT-09}

\par 難易度 2（校準初稿）、優先度未評価。

\section{社会的責任・安全の標準}

\subsection{平成23年度 問10}
\label{q:h23-reg-apt-10}

\noindent\textbf{原典：} sources/H23/aptitude.pdf、PDF物理11頁。

\textbf{公式正答：} 3（single）。

\par\textit{原典本文の校正と独自解説は未作成。分類のみの登録項目。}

\subsection{平成29年度 問11}
\label{q:h29-reg-apt-11}

\noindent\textbf{原典：} sources/H29/aptitude.pdf、PDF物理11頁。

\textbf{公式正答：} 5（single）。

\par\textit{原典本文の校正と独自解説は未作成。分類のみの登録項目。}

\subsection{令和4年度 問3}
\label{q:r04-reg-apt-03}

\noindent\textbf{原典：} sources/R04/aptitude.pdf、PDF物理2頁。

\textbf{公式正答：} 3（single）。

\par\textit{原典本文の校正と独自解説は未作成。分類のみの登録項目。}

\subsection{令和7年度 問5}
\label{q:r07-reg-apt-05}

\noindent\textbf{原典：} sources/R07/aptitude.pdf、PDF物理5頁。

\textbf{公式正答：} 1（single）。

\par\textit{原典本文の校正と独自解説は未作成。分類のみの登録項目。}
